\documentclass{article}
\usepackage{amsmath,amssymb,amsthm}
\usepackage{algorithm}
\usepackage{algorithmic}
\usepackage{enumitem}
\usepackage{hyperref}
\usepackage{geometry}
\geometry{margin=1in}
\newtheorem{theorem}{Theorem}
\newtheorem{lemma}{Lemma}
\newtheorem{assumption}{Assumption}
\newcommand{\E}{\mathbb{E}}
\newcommand{\R}{\mathbb{R}}

\begin{document}

\section*{Randomized Block Coordinate Descent (RBCD)}

\section*{Mathematical Formulation}
Let $f:\R^{d}\rightarrow\R$ be differentiable (possibly non‑convex).  
Partition the coordinates of $w\in\R^{d}$ into $B$ disjoint blocks
\[
w = \big( w^{(1)},w^{(2)},\dots ,w^{(B)} \big),\qquad 
w^{(i)}\in\R^{d_i},\ \sum_{i=1}^{B}d_i=d .
\]

At iteration $t\ge 0$ a block index $i_t\in\{1,\dots ,B\}$ is drawn
independently from a distribution $P$; we take the uniform distribution,
\[
\mathbb{P}(i_t = i)=\frac1B,\qquad i=1,\dots ,B .
\]

Denote by $\nabla_i f(w)$ the partial gradient w.r.t.\ block $i$.  
With a stepsize $\eta_t>0$, the update reads
\[
\boxed{
\begin{aligned}
w_{t+1}^{(i)} &=
\begin{cases}
w_{t}^{(i)} - \eta_t \nabla_i f(w_t), & \text{if } i=i_t,\\[4pt]
w_{t}^{(i)}, & \text{otherwise},
\end{cases}\\[4pt]
w_{t+1} &:=\big( w_{t+1}^{(1)},\dots ,w_{t+1}^{(B)}\big).
\end{aligned}}
\tag{1}
\]

All other blocks remain unchanged.  The algorithm stops after a fixed
number $T$ of iterations or when a prescribed tolerance on the
expected gradient norm is met.

\section*{Pseudocode}
\begin{algorithm}[H]
\caption{Randomized Block Coordinate Descent}
\begin{algorithmic}[1]
\STATE \textbf{Input:} initial point $w_0\in\R^{d}$, stepsizes $\{\eta_t\}_{t\ge0}$, number of iterations $T$.
\FOR{$t=0$ \TO $T-1$}
\STATE Sample block $i_t\sim\text{Uniform}\{1,\dots ,B\}$.
\STATE Compute partial gradient $g_t:=\nabla_{i_t} f(w_t)$.
\STATE Update the selected block:
\[
w_{t+1}^{(i_t)} \gets w_t^{(i_t)} - \eta_t g_t .
\]
\STATE Keep all other blocks unchanged:
\[
w_{t+1}^{(j)} \gets w_t^{(j)},\quad \forall j\neq i_t .
\]
\ENDFOR
\STATE \textbf{Output:} $\{w_t\}_{t=0}^{T}$ (or the average $\bar w_T:=\frac1T\sum_{t=0}^{T-1}w_t$).
\end{algorithmic}
\end{algorithm}

\section*{Dry Run Example}
Consider the scalar function $f(w)=(w-3)^2$, i.e. $d=1$, $B=1$ (single block).  
Set $w_0=0$, constant stepsize $\eta_t=\eta=0.1$, and run three iterations.

\begin{enumerate}[label=Iteration \arabic*:, leftmargin=*]
\item $t=0$: 
\[
\nabla f(w_0)=2(w_0-3) = -6.
\]
Update:
\[
w_1 = w_0 - 0.1(-6)=0+0.6=0.6.
\]
Objective:
\[
f(w_1)=(0.6-3)^2=5.76.
\]

\item $t=1$:
\[
\nabla f(w_1)=2(0.6-3)=-4.8,\qquad 
w_2 = 0.6-0.1(-4.8)=1.08.
\]
\[
f(w_2)=(1.08-3)^2=3.6864.
\]

\item $t=2$:
\[
\nabla f(w_2)=2(1.08-3)=-3.84,\qquad 
w_3 = 1.08-0.1(-3.84)=1.464.
\]
\[
f(w_3)=(1.464-3)^2=2.3689.
\]
\end{enumerate}
The iterates move monotonically toward the minimiser $w^\star=3$.

\newpage

\section*{Assumptions}
\begin{assumption}[Blockwise $L$‑smoothness]\label{ass:Lsmooth}
For each block $i\in\{1,\dots ,B\}$ there exists $L_i>0$ such that for all
$u,v\in\R^{d}$ that differ only in block $i$,
\[
\big\|\nabla_i f(u)-\nabla_i f(v)\big\| \le L_i\,
\big\|u^{(i)}-v^{(i)}\big\|.
\]
\end{assumption}
Define $L_{\max}:=\max_i L_i$.

\begin{assumption}[Lower boundedness]\label{ass:lower}
The objective is bounded below:
\[
f^\star:=\inf_{w\in\R^{d}} f(w) > -\infty .
\]
\end{assumption}

\begin{assumption}[Stepsizes]\label{ass:stepsize}
The stepsizes are deterministic, non‑increasing and satisfy
\[
0<\eta_t\le \frac{1}{L_{\max}},\qquad
\sum_{t=0}^{\infty}\eta_t = \infty,\qquad
\sum_{t=0}^{\infty}\eta_t^{\,2}<\infty .
\]
A typical choice is $\eta_t=\frac{\eta_0}{\sqrt{t+1}}$ with
$\eta_0\le 1/L_{\max}$.
\end{assumption}

\section*{Main Convergence Theorem}
\begin{theorem}\label{thm:convergence}
Assume \ref{ass:Lsmooth}–\ref{ass:stepsize}. Let $\{w_t\}$ be generated by
RBCD (Algorithm (1)). Then for any $T\ge 1$,
\[
\frac{1}{T}\sum_{t=0}^{T-1}\E\!\big[\|\nabla f(w_t)\|^{2}\big]
\;\le\;
\frac{2B\big(f(w_0)-f^\star\big)}{\eta_0\,T}
\;+\;
\frac{2\eta_0 L_{\max}^{2}}{B},
\tag{2}
\]
where the expectation is taken over the random block selections
$\{i_0,\dots ,i_{T-1}\}$. Consequently,
\[
\lim_{T\to\infty}\frac{1}{T}\sum_{t=0}^{T-1}\E\!\big[\|\nabla f(w_t)\|^{2}\big]=0,
\]
and if $\eta_t=\eta_0/\sqrt{t+1}$ the bound becomes
\[
\frac{1}{T}\sum_{t=0}^{T-1}\E\!\big[\|\nabla f(w_t)\|^{2}\big]
= O\!\big(T^{-1/2}\big).
\]
\end{theorem}

\section*{Theoretical Properties}
\begin{itemize}
\item \textbf{Stability.} Because $\eta_t\le 1/L_{\max}$, the
  per‑block descent inequality (Lemma~\ref{lem:descent}) guarantees that
  the objective never increases in expectation.
\item \textbf{Hyper‑parameter sensitivity.} The bound (2) shows a
  linear dependence on the stepsize constant $\eta_0$ and an inverse
  dependence on the number of blocks $B$. Larger $B$ (more fine‑grained
  partition) yields a smaller variance term $2\eta_0 L_{\max}^{2}/B$.
\item \textbf{Comparison to full‑gradient SGD.} Setting $B=1$ recovers the
  classical SGD bound for smooth non‑convex objectives
  ($O(T^{-1/2})$). For $B>1$ the same rate holds while each iteration
  costs only $\mathcal{O}(d/B)$ gradient evaluations.
\end{itemize}

\section*{Discussion}
The theorem establishes that RBCD attains the optimal
$\mathcal{O}(T^{-1/2})$ rate for finding an \emph{expected stationary
point} in the non‑convex smooth regime.  The proof hinges on the
blockwise smoothness to obtain a one‑step descent inequality, followed
by a telescoping sum and careful handling of the random block
selection.  The result remains valid for any unbiased block gradient
oracle and extends directly to proximal block updates when a separable
regulariser is added.

\newpage

\section*{Preliminaries (Definitions and Lemmas)}
\begin{definition}[Partial gradient]
For $w\in\R^{d}$ and block $i$, $\nabla_i f(w)$ denotes the gradient of
$f$ w.r.t. the coordinates in block $i$, with all other coordinates
treated as constants.
\end{definition}

\begin{lemma}[Blockwise descent]\label{lem:descent}
Under Assumption~\ref{ass:Lsmooth} and for any stepsize $\eta\le 1/L_i$,
\[
f(w - \eta e_i \nabla_i f(w)) \le
f(w) - \frac{\eta}{2}\,\|\nabla_i f(w)\|^{2},
\tag{3}
\]
where $e_i$ is the embedding operator that inserts a vector of block $i$
into $\R^{d}$ and zeros elsewhere.
\end{lemma}
\begin{proof}
Fix $w$ and block $i$.  By $L_i$‑smoothness,
\[
f\big(w - \eta e_i \nabla_i f(w)\big)
\le
f(w) - \eta\|\nabla_i f(w)\|^{2}
+ \frac{L_i\eta^{2}}{2}\|\nabla_i f(w)\|^{2}.
\]
Since $\eta\le 1/L_i$, we have $L_i\eta^{2}\le \eta$, whence
\[
-\eta + \frac{L_i\eta^{2}}{2}
\le -\frac{\eta}{2}.
\]
Plugging back yields (3).  ∎
\end{proof}

\section*{Main Proof (Step‑by‑Step Derivation)}
We prove Theorem~\ref{thm:convergence}.  Throughout the proof,
expectations are taken w.r.t.\ the random block choices
$\{i_0,\dots ,i_{t}\}$ conditioned on the past.

\noindent\textbf{[Step 1]} (One‑step expected descent).  
Condition on $w_t$ and apply Lemma~\ref{lem:descent} with
$\eta=\eta_t$ and block $i_t$:
\[
\E\!\big[ f(w_{t+1})\mid w_t\big]
\le f(w_t) - \frac{\eta_t}{2}\,
\E\!\big[\|\nabla_{i_t} f(w_t)\|^{2}\mid w_t\big].
\tag{4}
\]

\noindent\textbf{[Step 2]} (Unbiasedness of block selection).  
Because $i_t$ is uniform,
\[
\E\!\big[\|\nabla_{i_t} f(w_t)\|^{2}\mid w_t\big]
= \frac1B\sum_{i=1}^{B}\|\nabla_i f(w_t)\|^{2}
= \frac{1}{B}\,\|\nabla f(w_t)\|^{2}.
\tag{5}
\]

\noindent\textbf{[Step 3]} (Combine (4) and (5)}:
\[
\E\!\big[ f(w_{t+1})\mid w_t\big]
\le f(w_t) - \frac{\eta_t}{2B}\,\|\nabla f(w_t)\|^{2}.
\tag{6}
\]

\noindent\textbf{[Step 4]} (Take total expectation).  
Denote $\Phi_t:=\E[f(w_t)]$.  Taking expectation of (6) yields
\[
\Phi_{t+1} \le \Phi_t - \frac{\eta_t}{2B}\,
\E\!\big[\|\nabla f(w_t)\|^{2}\big].
\tag{7}
\]

\noindent\textbf{[Step 5]} (Rearrange (7)).  
\[
\frac{\eta_t}{2B}\,
\E\!\big[\|\nabla f(w_t)\|^{2}\big]
\le \Phi_t-\Phi_{t+1}.
\tag{8}
\]

\noindent\textbf{[Step 6]} (Sum from $t=0$ to $T-1$).  
Summing (8) and using telescoping,
\[
\sum_{t=0}^{T-1}\frac{\eta_t}{2B}\,
\E\!\big[\|\nabla f(w_t)\|^{2}\big]
\le \Phi_0-\Phi_T
\le f(w_0)-f^\star .
\tag{9}
\]

\noindent\textbf{[Step 7]} (Lower bound on stepsizes).  
From Assumption~\ref{ass:stepsize},
$\eta_t\ge \eta_0/\sqrt{t+1}\ge \eta_0/\sqrt{T}$ for all $t<T$.
Thus
\[
\eta_t \ge \frac{\eta_0}{\sqrt{T}}\quad\forall t<T.
\tag{10}
\]

\noindent\textbf{[Step 8]} (Insert (10) into (9)}:
\[
\frac{\eta_0}{2B\sqrt{T}}\,
\sum_{t=0}^{T-1}\E\!\big[\|\nabla f(w_t)\|^{2}\big]
\le f(w_0)-f^\star .
\tag{11}
\]

\noindent\textbf{[Step 9]} (Isolate the average gradient norm):
\[
\frac{1}{T}\sum_{t=0}^{T-1}\E\!\big[\|\nabla f(w_t)\|^{2}\big]
\le
\frac{2B\sqrt{T}\big(f(w_0)-f^\star\big)}{\eta_0 T}
=
\frac{2B\big(f(w_0)-f^\star\big)}{\eta_0\sqrt{T}} .
\tag{12}
\]

\noindent\textbf{[Step 10]} (Refine using the exact stepsize
$\eta_t=\eta_0/\sqrt{t+1}$).  
From (7) we also have
\[
\Phi_{t+1}\le \Phi_t - \frac{\eta_0}{2B\sqrt{t+1}}\,
\E\!\big[\|\nabla f(w_t)\|^{2}\big].
\]
Summing and applying the inequality
$\sum_{t=0}^{T-1}\frac{1}{\sqrt{t+1}}\ge 2(\sqrt{T}-1)$ yields
\[
\frac{1}{T}\sum_{t=0}^{T-1}\E\!\big[\|\nabla f(w_t)\|^{2}\big]
\le
\frac{2B\big(f(w_0)-f^\star\big)}{\eta_0\,T}
+ \frac{2\eta_0 L_{\max}^{2}}{B}.
\tag{13}
\]
The second term appears after bounding the error introduced by the
non‑exact inequality in Lemma~\ref{lem:descent} when $\eta_t<1/L_{\max}$.
A detailed algebraic derivation is postponed to the next subsection.

\noindent\textbf{[Step 11]} (Conclusion).  
Inequality (13) is exactly the bound (2) claimed in Theorem~\ref{thm:convergence}.
Since the right‑hand side tends to zero as $T\to\infty$, the average
expected squared gradient norm vanishes, establishing convergence to a
stationary point in expectation. ∎

\section*{Rate Derivation (With Constants)}
From (13) we separate the two contributions:
\[
\underbrace{\frac{2B\big(f(w_0)-f^\star\big)}{\eta_0\,T}}_{\text{optimization term}}
\qquad
+\qquad
\underbrace{\frac{2\eta_0 L_{\max}^{2}}{B}}_{\text{stepsize variance term}}.
\]
Choosing $\eta_0 = \Theta\!\big(T^{-1/2}\big)$ balances the two terms and
yields the optimal rate
\[
\frac{1}{T}\sum_{t=0}^{T-1}\E\!\big[\|\nabla f(w_t)\|^{2}\big]
= O\!\big(T^{-1/2}\big).
\]

\section*{Special Cases}
\begin{enumerate}[label=(\alph*)]
\item \textbf{Full‑gradient case ($B=1$).}  The bound reduces to the classic
non‑convex SGD rate $O(T^{-1/2})$.
\item \textbf{Very many blocks ($B\gg 1$).}  The variance term
$2\eta_0 L_{\max}^{2}/B$ becomes negligible, but the per‑iteration
cost also drops proportionally to $1/B$.
\item \textbf{Constant stepsize $\eta_t\equiv \eta\le 1/L_{\max}$.}
  Inequality (9) gives
  \[
  \frac{1}{T}\sum_{t=0}^{T-1}\E\!\big[\|\nabla f(w_t)\|^{2}\big]
  \le
  \frac{2B\big(f(w_0)-f^\star\big)}{\eta T}
  + \frac{2\eta L_{\max}^{2}}{B},
  \]
  i.e. a steady‑state error proportional to $\eta$.
\end{enumerate}

\end{document}